\uuid{t0x1}
\exo7id{3824}
\titre{exo7 3824}
\auteur{quercia}
\organisation{exo7}
\datecreate{2010-03-11}
\isIndication{false}
\isCorrection{false}
\chapitre{Espace euclidien, espace normé}
\sousChapitre{Espaces vectoriels hermitiens}
\cpge{mpsi-mp2i.espaces-vectoriels.a, pcsi.espaces-vectoriels.a, ptsi.geometrie-espace.c, mpsi-mp2i.espaces-prehilbertiens.a, pcsi.espaces-prehilbertiens.a, ptsi.geometrie-plan.b, mpsi-mp2i.espaces-prehilbertiens.b, pcsi.espaces-prehilbertiens.b, ptsi.geometrie-espace.b, mp-mpi.topologie-espaces-normes.a, pc.espaces-vectoriels-normes.a, psi.espaces-vectoriels-normes.a, mp-mpi.series-numeriques-vectorielles.a, pt.series-numeriques.b}
\module{Algèbre}
\niveau{L2}
\difficulte{}

\contenu{
\texte{
Soit $E$ l'ensemble des suites $(u_n)_{n\in\N}$ à termes réels telles que
la série $\sum u_n^2$ converge.

Pour $u,v \in E$, on pose : $(u\mid v) = \sum_{n=0}^{\infty} u_nv_n$.
}
\begin{enumerate}
    \item \question{Montrer que $E$ est un espace vectoriel sur $\R$.}
    \item \question{Montrer que $(u\mid v)$ existe.}
    \item \question{Montrer qu'on définit ainsi un produit scalaire sur $E$.}
    \item \question{Montrer que $E$, muni de la norme associée, est complet.}
\end{enumerate}
}
